\clearpage
\section{PROJETO CONCEITUAL}

O projeto conceitual busca meios de realizar as funções necessárias ao atendimento das especificações-meta. A passagem é de requisito para função, de função para princípio e, depois, para uma concepção verificável. A modelagem representa o que o produto deve fazer antes de fixar materiais, dimensões ou componentes \citep{back2008,rozenfeld2006}.

A análise funcional compreende sete parciais, 21 elementares da M1 e sete complementares de retorno, manutenção e fim de vida. O FAST relaciona finalidade e meios, enquanto três estruturas explicitam alternativas para a proteção e a conferência. A matriz de Pugh seleciona M1, proteção direta da peça após preparo no distribuidor, com total ponderado +7. A decisão compara as capacidades, as interfaces e as contrapartidas das três modelagens \citep{equipeconceitual2026,planilhapsp62026}.

\subsection{Modelagem Funcional}

O fluxo de aplicação compreende as etapas de distribuidor, transporte, descarga e movimentação na obra. A louça é o material protegido; o produto em desenvolvimento é o dispositivo retornável. A delimitação de porte médio considera cubas e corpos individuais de bacias, uma peça por unidade. Kits montados e monoblocos não foram incluídos automaticamente nessa faixa.

A situação atual, V0, é a louça acondicionada na embalagem do fabricante. A observação de uma obra residencial e as duas entrevistas sustentam a investigação da proteção e da conferência, mas não quantificam a taxa de quebra nem demonstram em que operação o dano ocorreu. Na modelagem, o papelão opaco é uma barreira ao caminho visual até a cerâmica. A solução precisa tratar a condição da peça e a segurança da retirada, além de transportar o volume.

O procedimento segue a atividade de modelar funcionalmente o produto: analisar as especificações-meta; identificar funções; estabelecer a função global e os fluxos; desenvolver estruturas alternativas; e selecionar uma estrutura para a geração de princípios. Funções usam verbo no infinitivo e substantivo. Qualidades e limites são associados aos efeitos e às interfaces que precisam satisfazê-los \citep{aulasmodelagem}.

\subsubsection{Função global e fluxos de material, energia e sinal}

A função global foi expressa como \textbf{Proteger louça}. O fluxo material principal é a peça que precisa atravessar a operação preservando sua integridade. Se ocorrer quebra, os fragmentos devem permanecer contidos. Choques, vibração, compressão e manuseio constituem solicitações mecânicas; parte da energia é dissipada e parte pode ser transmitida à peça. Os sinais abrangem a condição observável, o fechamento e a identidade da unidade.

\begin{figure}[H]
\centering\begin{tikzpicture}[x=.995mm,y=1mm]
\tikzset{
 fn/.style={draw=cinzamedio,fill=white,align=center,font=\footnotesize,inner xsep=1mm,inner ysep=1mm,text width=22mm,minimum height=13mm,execute at begin node=\setstretch{1}},
 fp/.style={fn,fill=cinzaclaro,minimum height=14mm},
 material/.style={-{Stealth[length=1.7mm]},draw=cinzaescuro,line width=.8pt},
 energia/.style={-{Stealth[length=1.5mm]},draw=cinzamedio,dashed,line width=.6pt},
 sinal/.style={-{Stealth[length=1.5mm]},draw=cinzamedio,dash pattern=on 3pt off 2pt on .6pt off 2pt,line width=.6pt},
 rel/.style={-{Stealth[length=1mm]},draw=cinzamedio,line width=.45pt},
 porta/.style={align=left,font=\footnotesize,inner sep=0,execute at begin node=\setstretch{1}},
 legenda/.style={font=\footnotesize,anchor=west,execute at begin node=\setstretch{1}}
}

\node[fp,text width=32mm,minimum height=67mm,font=\normalsize] (fg) at (80,-27) {\textbf{Proteger louça}};
\node[porta,anchor=south west,text width=58mm] at (0,1) {\textbf{Material:} louça a proteger};
\node[porta,anchor=south west,text width=58mm] at (0,-17) {\textbf{Energia:} choque, vibração e compressão};
\node[porta,anchor=south west,text width=58mm] at (0,-35) {\textbf{Energia:} trabalho mecânico do operador};
\node[porta,anchor=south west,text width=58mm] at (0,-53) {\textbf{Sinal:} condição da louça a revelar};
\draw[material] (0,0)--(62,0);\draw[energia] (0,-18)--(62,-18);\draw[energia] (0,-36)--(62,-36);\draw[sinal] (0,-54)--(62,-54);
\node[porta,anchor=south west,text width=58.5mm] at (101,1) {\textbf{Material:} louça íntegra};
\node[porta,anchor=south west,text width=58.5mm] at (101,-17) {\textbf{Material:} fragmentos retidos, se houver quebra};
\node[porta,anchor=south west,text width=58.5mm] at (101,-35) {\textbf{Energia:} energia dissipada e parcela transmitida};
\node[porta,anchor=south west,text width=58.5mm] at (101,-53) {\textbf{Sinal:} condição, fechamento e identidade};
\draw[material] (98,0)--(160,0);\draw[material] (98,-18)--(160,-18);\draw[energia] (98,-36)--(160,-36);\draw[sinal] (98,-54)--(160,-54);
\draw[material] (0,-73)--(14,-73);\node[legenda] at (17,-73) {Material};
\draw[energia] (55,-73)--(69,-73);\node[legenda] at (72,-73) {Energia};
\draw[sinal] (110,-73)--(124,-73);\node[legenda] at (127,-73) {Sinal};
\end{tikzpicture}

\caption{Função global do dispositivo e fluxos de material, energia e sinal}
\label{fig:funcao-global}
\fonte{Elaboração própria.}
\end{figure}

O trabalho mecânico do operador ajusta, trava, movimenta e abre o conjunto. A informação de condição pode ser visual e passiva; a presença de sinal no diagrama não exige eletrônica. Os estados de saída da Figura~\ref{fig:funcao-global} são efeitos pretendidos. Sua representação não comprova proteção, retenção ou inspeção na área crítica.

\subsubsection{Relação entre requisitos, funções e restrições}
\label{sec:requisitos-funcoes}

A análise das especificações-meta distingue o efeito que o produto precisa realizar, o desempenho exigido e as condições de uso. As funções são escritas por verbo no infinitivo e substantivo. A redação sintetiza uma ação do sistema técnico sobre material, energia ou sinal; a meta determina em que condição essa ação deverá ser verificada \citep{back2008,aulasmodelagem}.

Dos 23 RP, quinze orientam a identificação de funções e oito foram tratados como parâmetros ou restrições: RP8, RP10, RP12, RP13, RP14, RP18, RP20 e RP23. Essa classificação não elimina requisitos nem impede desdobrar capacidades auxiliares. A facilidade de uso, por exemplo, qualifica a interação e orienta Indicar fechamento; a modularidade restringe as interfaces e exige viabilizar a reposição. Uma função pode atender vários requisitos e um requisito pode depender de várias funções.

\begin{table}[H]
\caption{Tratamento dos requisitos não funcionais na modelagem}
\label{tab:requisito-funcao}
\centering\footnotesize
\begin{tabularx}{\textwidth}{L{1.0cm}L{3.0cm}Y}
\toprule\textbf{RP}&\textbf{Atributo ou restrição}&\textbf{Aplicação às funções e à solução}\\\midrule
RP8&Compatibilidade com o veículo&Limita envelope, apoio e amarração no palete e no veículo. Não cria um bloco de planejamento de entregas.\\
RP10&Uso intuitivo&Qualifica a interface de uso; E17, Indicar fechamento, contribui para orientar a ação.\\
RP12&Uniformidade&Qualifica a resposta de E04/E05; comparar o pico entre unidades nas mesmas condições.\\
RP13&Resistência à umidade&Qualifica a capacidade do caminho resistente de E06 após condicionamento.\\
RP14&Volume externo&Restringe o conjunto carregado por família. C01, Recolher estrutura, trata do estado vazio no retorno.\\
RP18&Vida em ciclos&Qualifica a permanência das capacidades de proteção, apoio, acesso e união ao longo do reuso.\\
RP20&Modularidade&Restringe arquitetura e interfaces. C04/C05 viabilizam reposição e união; comprovar a troca independente.\\
RP23&Custo por ciclo&Restringe todas as funções e soluções. Comparar os custos na mesma fronteira operacional.\\
\bottomrule\end{tabularx}
\fonte{Análise dos 23 requisitos de projeto e seus efeitos funcionais.}
\end{table}

Os limites preservam a especificidade do produto. E04, Absorver impacto, deverá contribuir para manter o pico na cerâmica em até 35\,g e abaixo da embalagem atual no mesmo ensaio. E03, Restringir movimento, está associado ao deslocamento residual de até 6\,mm. E16, Expor superfície, deverá permitir conferir pelo menos 60\% da área crítica com a contenção fechada. E11, Ajustar apoios, está vinculado aos cursos de 300\,mm em A e 150\,mm em B, sem troca de componentes. A presença dessas funções no desenho não demonstra atendimento das metas.

\subsubsection{Funções parciais e elementares}

A função global foi desdobrada em sete parciais. Imobilizar louça, Atenuar esforços e Conter fragmentos representam as capacidades básicas de proteção; Acomodar louça, Movimentar conjunto, Informar condição e Liberar louça organizam a interação com a peça e com o usuário. Cada parcial foi detalhada em três elementares na M1, conforme a Tabela~\ref{tab:catalogo-funcoes}.

\begin{table}[H]
\caption{Funções parciais e 21 elementares da M1}
\label{tab:catalogo-funcoes}
\centering\footnotesize
\begin{tabularx}{\textwidth}{L{0.7cm}L{3.0cm}L{2.0cm}Y}
\toprule\textbf{ID}&\textbf{Função parcial}&\textbf{Classe}&\textbf{Funções elementares}\\\midrule
\linhasquadro{quadros/parciais-elementares}
\bottomrule\end{tabularx}
\fonte{Elaboração própria: catálogo funcional da M1.}
\end{table}

Travar apoios fixa a regulagem; Distribuir pressão evita concentrar a carga de contato; Restringir movimento conserva a posição relativa da louça. Esses efeitos precisam ser compatíveis: aumentar a restrição sem controlar a pressão pode danificar a cerâmica. Absorver impacto e Isolar vibração tratam das solicitações dinâmicas; Desviar compressão exige um caminho resistente que preserve a peça e não anule a proteção dinâmica.

Expor superfície significa disponibilizar um caminho visual até a região crítica da própria cerâmica, com a contenção fechada. Indicar fechamento comunica o estado e a orientação de uso. Exibir identidade apresenta o identificador único do dispositivo. As três funções tratam de sinais diferentes. Um número de série não orienta, por si só, a montagem; um aviso de choque não torna a trinca visível.

As funções técnicas atuam na sustentação, na restrição do movimento, na transmissão de esforços e na contenção. As funções de transformação tratam da resposta ao impacto e à vibração. As interativas abrangem interfaces ergonômicas, leitura do uso e comunicação: Receber força disponibiliza a pega; Abrir barreira e Destravar apoios viabilizam acesso; Indicar fechamento e Exibir identidade comunicam estado e significado. Essa classificação caracteriza os efeitos e as interfaces, sem substituir a distinção entre função e requisito de qualidade \citep{aulasmodelagem}.

A equipe também relacionou as funções a percepções pretendidas: confiança na proteção, segurança na abertura e autonomia de uso. Essas percepções deverão ser confrontadas com usuários. Elas não comprovam aceleração, retenção ou usabilidade e não foram usadas como funções técnicas adicionais para aumentar a pontuação.

Os IDs P1--P7 e E01--E21 identificam os mesmos efeitos nos diferentes níveis de representação. Nas estruturas sobre o volume embalado, os nomes Receber volume, Posicionar volume, Envolver volume e Remover volume explicitam a mudança de material, mantendo o ID do efeito equivalente. Os IDs não representam peças nem módulos definidos.

\subsubsection{Funções complementares do ciclo de vida}

O caráter retornável exige capacidades além da proteção da louça durante uma entrega. Sete complementares sustentam retorno, manutenção e fim de vida. Elas são comuns às três alternativas e permanecem necessárias após a seleção funcional.

\begin{table}[H]
\caption{Funções complementares comuns às modelagens}
\label{tab:complementares}
\centering\footnotesize
\begin{tabularx}{\textwidth}{L{0.8cm}L{3.2cm}YL{2.0cm}}
\toprule\textbf{ID}&\textbf{Função}&\textbf{Efeito ou capacidade}&\textbf{RP}\\\midrule
\linhasquadro{quadros/funcoes-complementares}
\bottomrule\end{tabularx}
\fonte{Catálogo funcional e vínculos com o ciclo de vida do dispositivo retornável.}
\end{table}

\begin{figure}[H]
\centering\begin{tikzpicture}[x=1mm,y=1mm]
\tikzset{
 fn/.style={draw=cinzamedio,fill=white,align=center,font=\footnotesize,inner xsep=1mm,inner ysep=1mm,text width=22mm,minimum height=13mm,execute at begin node=\setstretch{1}},
 fp/.style={fn,fill=cinzaclaro,minimum height=14mm},
 material/.style={-{Stealth[length=1.7mm]},draw=cinzaescuro,line width=.8pt},
 energia/.style={-{Stealth[length=1.5mm]},draw=cinzamedio,dashed,line width=.6pt},
 sinal/.style={-{Stealth[length=1.5mm]},draw=cinzamedio,dash pattern=on 3pt off 2pt on .6pt off 2pt,line width=.6pt},
 rel/.style={-{Stealth[length=1mm]},draw=cinzamedio,line width=.45pt},
 porta/.style={align=left,font=\footnotesize,inner sep=0,execute at begin node=\setstretch{1}},
 legenda/.style={font=\footnotesize,anchor=west,execute at begin node=\setstretch{1}}
}

\node[fp,text width=44mm] at (25,0) {\textbf{Preparar retorno}};
\node[fn,text width=44mm] at (25,-22) {{\scriptsize C01}\\\textbf{Recolher estrutura}};
\node[fn,text width=44mm] at (25,-42) {{\scriptsize C02}\\\textbf{Encaixar unidades}};
\node[fp,text width=44mm] at (80,0) {\textbf{Recondicionar conjunto}};
\node[fn,text width=44mm] at (80,-22) {{\scriptsize C03}\\\textbf{Expor interfaces}};
\node[fn,text width=44mm] at (80,-42) {{\scriptsize C04}\\\textbf{Substituir módulos}};
\node[fn,text width=44mm] at (80,-62) {{\scriptsize C05}\\\textbf{Unir partes}};
\node[fp,text width=44mm] at (135,0) {\textbf{Desmontar conjunto}};
\node[fn,text width=44mm] at (135,-22) {{\scriptsize C06}\\\textbf{Soltar uniões}};
\node[fn,text width=44mm] at (135,-42) {{\scriptsize C07}\\\textbf{Separar materiais}};
\end{tikzpicture}

\caption{Organização das capacidades de retorno, manutenção e fim de vida}
\label{fig:ciclo-vida-funcional}
\fonte{Elaboração própria a partir das funções complementares do ciclo de vida.}
\end{figure}

O operador limpa, inspeciona, ajusta e troca partes. O produto viabiliza essas operações por acessos, interfaces e uniões liberáveis. Substituir módulos expressa essa capacidade de reposição; sua realização depende de demonstrar interfaces e troca independente. A modelagem não atribui ao dispositivo a gestão de estoque, de rotas ou de entregas.

Recolher estrutura atua sobre o conjunto vazio; não comprova o volume carregado do RP14. Encaixar unidades apoia o arranjo entre dispositivos, mas não demonstra compressão nem estabilidade. Separar materiais viabiliza frações ao fim de vida, sem comprovar automaticamente a fração reciclável ou a cadeia regional.


\subsubsection{Aplicação do método FAST}

O FAST relaciona a finalidade aos meios de realizá-la. A leitura \textit{como?} parte de Proteger louça para as parciais e suas elementares; a leitura \textit{por quê?} retorna à finalidade. Reduzir avarias e evitar cortes são finalidades superiores que orientam a modelagem. O dispositivo oferece as capacidades físicas para contribuir com esses resultados, sem controlar todas as causas do processo \citep{back2008,aulasmodelagem}.

Na parcial Imobilizar louça, por exemplo, a restrição depende de travar os apoios e distribuir a pressão sem permitir movimento danoso. Em Informar condição, expor a superfície, indicar fechamento e exibir identidade atendem necessidades informacionais distintas. As três elementares de cada linha estão no mesmo nível de desdobramento; não representam três níveis sucessivos ou uma sequência cronológica de operações.

\begin{landscape}
\begin{figure}[H]
\centering\begin{tikzpicture}[x=1mm,y=1mm]
\tikzset{
 fn/.style={draw=cinzamedio,fill=white,align=center,font=\footnotesize,inner xsep=1mm,inner ysep=1mm,text width=22mm,minimum height=13mm,execute at begin node=\setstretch{1}},
 fp/.style={fn,fill=cinzaclaro,minimum height=14mm},
 material/.style={-{Stealth[length=1.7mm]},draw=cinzaescuro,line width=.8pt},
 energia/.style={-{Stealth[length=1.5mm]},draw=cinzamedio,dashed,line width=.6pt},
 sinal/.style={-{Stealth[length=1.5mm]},draw=cinzamedio,dash pattern=on 3pt off 2pt on .6pt off 2pt,line width=.6pt},
 rel/.style={-{Stealth[length=1mm]},draw=cinzamedio,line width=.45pt},
 porta/.style={align=left,font=\footnotesize,inner sep=0,execute at begin node=\setstretch{1}},
 legenda/.style={font=\footnotesize,anchor=west,execute at begin node=\setstretch{1}}
}

\node[porta,anchor=west] at (1,17) {\textbf{Por quê?} $\longleftarrow$};
\node[porta,anchor=east] at (244,17) {$\longrightarrow$ \textbf{Como?}};
\node[fp,text width=32mm,minimum height=20mm] (fg) at (27,-51) {\textbf{Proteger louça}};
\node[fp,text width=40mm] (p0) at (88,0) {{\scriptsize P1}\\\textbf{Imobilizar louça}};
\draw[rel] (fg.east) -- (56,-51) |- (p0.west);
\node[fn,text width=35mm] (e00) at (137,0) {{\scriptsize E01}\\\textbf{Travar apoios}};
\node[fn,text width=35mm] (e01) at (180,0) {{\scriptsize E02}\\\textbf{Distribuir pressão}};
\node[fn,text width=35mm] (e02) at (223,0) {{\scriptsize E03}\\\textbf{Restringir movimento}};
\draw[draw=cinzamedio,line width=.45pt] (p0.east) -- (114,0) -- (114,9) -- (223,9);
\draw[rel] (137,9) -- (e00.north);
\draw[rel] (180,9) -- (e01.north);
\draw[rel] (223,9) -- (e02.north);
\node[fp,text width=40mm] (p1) at (88,-17) {{\scriptsize P2}\\\textbf{Atenuar esforços}};
\draw[rel] (fg.east) -- (56,-51) |- (p1.west);
\node[fn,text width=35mm] (e10) at (137,-17) {{\scriptsize E04}\\\textbf{Absorver impacto}};
\node[fn,text width=35mm] (e11) at (180,-17) {{\scriptsize E05}\\\textbf{Isolar vibração}};
\node[fn,text width=35mm] (e12) at (223,-17) {{\scriptsize E06}\\\textbf{Desviar compressão}};
\draw[draw=cinzamedio,line width=.45pt] (p1.east) -- (114,-17) -- (114,-8) -- (223,-8);
\draw[rel] (137,-8) -- (e10.north);
\draw[rel] (180,-8) -- (e11.north);
\draw[rel] (223,-8) -- (e12.north);
\node[fp,text width=40mm] (p2) at (88,-34) {{\scriptsize P3}\\\textbf{Conter fragmentos}};
\draw[rel] (fg.east) -- (56,-51) |- (p2.west);
\node[fn,text width=35mm] (e20) at (137,-34) {{\scriptsize E07}\\\textbf{Envolver louça}};
\node[fn,text width=35mm] (e21) at (180,-34) {{\scriptsize E08}\\\textbf{Reter fragmentos}};
\node[fn,text width=35mm] (e22) at (223,-34) {{\scriptsize E09}\\\textbf{Vedar juntas}};
\draw[draw=cinzamedio,line width=.45pt] (p2.east) -- (114,-34) -- (114,-25) -- (223,-25);
\draw[rel] (137,-25) -- (e20.north);
\draw[rel] (180,-25) -- (e21.north);
\draw[rel] (223,-25) -- (e22.north);
\node[fp,text width=40mm] (p3) at (88,-51) {{\scriptsize P4}\\\textbf{Acomodar louça}};
\draw[rel] (fg.east) -- (56,-51) |- (p3.west);
\node[fn,text width=35mm] (e30) at (137,-51) {{\scriptsize E10}\\\textbf{Receber louça}};
\node[fn,text width=35mm] (e31) at (180,-51) {{\scriptsize E11}\\\textbf{Ajustar apoios}};
\node[fn,text width=35mm] (e32) at (223,-51) {{\scriptsize E12}\\\textbf{Posicionar louça}};
\draw[draw=cinzamedio,line width=.45pt] (p3.east) -- (114,-51) -- (114,-42) -- (223,-42);
\draw[rel] (137,-42) -- (e30.north);
\draw[rel] (180,-42) -- (e31.north);
\draw[rel] (223,-42) -- (e32.north);
\node[fp,text width=40mm] (p4) at (88,-68) {{\scriptsize P5}\\\textbf{Movimentar conjunto}};
\draw[rel] (fg.east) -- (56,-51) |- (p4.west);
\node[fn,text width=35mm] (e40) at (137,-68) {{\scriptsize E13}\\\textbf{Receber força}};
\node[fn,text width=35mm] (e41) at (180,-68) {{\scriptsize E14}\\\textbf{Transmitir força}};
\node[fn,text width=35mm] (e42) at (223,-68) {{\scriptsize E15}\\\textbf{Estabilizar conjunto}};
\draw[draw=cinzamedio,line width=.45pt] (p4.east) -- (114,-68) -- (114,-59) -- (223,-59);
\draw[rel] (137,-59) -- (e40.north);
\draw[rel] (180,-59) -- (e41.north);
\draw[rel] (223,-59) -- (e42.north);
\node[fp,text width=40mm] (p5) at (88,-85) {{\scriptsize P6}\\\textbf{Informar condição}};
\draw[rel] (fg.east) -- (56,-51) |- (p5.west);
\node[fn,text width=35mm] (e50) at (137,-85) {{\scriptsize E16}\\\textbf{Expor superfície}};
\node[fn,text width=35mm] (e51) at (180,-85) {{\scriptsize E17}\\\textbf{Indicar fechamento}};
\node[fn,text width=35mm] (e52) at (223,-85) {{\scriptsize E18}\\\textbf{Exibir identidade}};
\draw[draw=cinzamedio,line width=.45pt] (p5.east) -- (114,-85) -- (114,-76) -- (223,-76);
\draw[rel] (137,-76) -- (e50.north);
\draw[rel] (180,-76) -- (e51.north);
\draw[rel] (223,-76) -- (e52.north);
\node[fp,text width=40mm] (p6) at (88,-102) {{\scriptsize P7}\\\textbf{Liberar louça}};
\draw[rel] (fg.east) -- (56,-51) |- (p6.west);
\node[fn,text width=35mm] (e60) at (137,-102) {{\scriptsize E19}\\\textbf{Destravar apoios}};
\node[fn,text width=35mm] (e61) at (180,-102) {{\scriptsize E20}\\\textbf{Abrir barreira}};
\node[fn,text width=35mm] (e62) at (223,-102) {{\scriptsize E21}\\\textbf{Remover louça}};
\draw[draw=cinzamedio,line width=.45pt] (p6.east) -- (114,-102) -- (114,-93) -- (223,-93);
\draw[rel] (137,-93) -- (e60.north);
\draw[rel] (180,-93) -- (e61.north);
\draw[rel] (223,-93) -- (e62.north);
\end{tikzpicture}

\caption{FAST da M1: função global, sete parciais e 21 elementares}
\label{fig:fast}
\fonte{Elaboração própria com base na modelagem funcional do dispositivo.}
\notaquadro{As complementares da Figura~\ref{fig:ciclo-vida-funcional} apoiam as três estruturas. As conexões do FAST representam finalidade e meios.}
\end{figure}
\end{landscape}

\subsubsection{Referências para os efeitos funcionais}

O estudo documental de embalagens e dispositivos existentes orientou a identificação de apoios, encaixes, amortecimento, caminho resistente e acesso. As referências International Paper, Delta/Masco e Rengo oferecem analogias para posicionamento e transmissão de esforços. As aberturas das referências Kohler/Vavra ajudam a discutir acesso visual, mas não demonstram inspeção da área crítica do mix considerado com a contenção fechada. Os dossiês de anterioridade registram documentos e limites dessas analogias \citep{equipeanterioridade2026}.

Contenção de fragmentos, ajuste ao mix, estabilidade, comunicação, identidade e reposição derivam das necessidades e dos requisitos de projeto. A anterioridade fornece referências de princípios; sua combinação é julgada pela adequação ao problema, às interfaces e às metas. A pesquisa documental sustenta a comparação tecnológica, sem constituir avaliação de patenteabilidade.

\subsubsection{Modelagem Funcional 1: proteção direta da peça descoberta}

Na M1, a embalagem original é retirada no distribuidor e a proteção atua diretamente sobre a louça. Acomodar louça organiza recebimento, ajuste e posicionamento; Imobilizar louça fixa os apoios e restringe o movimento. Atenuar esforços trata de impacto, vibração e compressão. A contenção e a abertura devem permitir retirar a peça sem perder a proteção contra fragmentos.

Essa fronteira retira o papelão opaco do caminho de observação e pode reduzir o envelope. Em contrapartida, exige preparo e transferência no distribuidor, regulagem para geometrias irregulares, controle da pressão nos contatos e manutenção das regiões de apoio. A estrutura da Figura~\ref{fig:modelagem-m1} preserva inspeção pela própria barreira do dispositivo; essa capacidade precisa ser demonstrada na área crítica.

\subsubsection{Modelagem Funcional 2: proteção do volume embalado}

Na M2, a embalagem original permanece e o dispositivo atua sobre o volume. As funções equivalentes passam a Receber volume, Posicionar volume, Envolver volume e Remover volume. A manutenção da caixa reduz a intervenção inicial, mas conserva sua folga interna e sua interferência no caminho de observação, conforme a hipótese usada na comparação.

A parcial Informar condição acrescenta A01, Abrir acesso, antes de Expor superfície. Isso torna explícita uma dificuldade da alternativa: o acesso através do papelão precisa ser compatibilizado com o requisito de conferir com o conjunto fechado. Um visor apenas externo não resolve a interferência da embalagem interna. A Figura~\ref{fig:modelagem-m2} mostra esse desdobramento adicional.

\subsubsection{Modelagem Funcional 3: volume embalado com monitoramento}

A M3 conserva a cadeia material e as funções da M2, mostradas na Figura~\ref{fig:modelagem-m3}, e acrescenta Monitorar impacto em paralelo. O ramo desdobra-se em Detectar aceleração, Comparar limiar, Registrar evento e Exibir alerta. Ele converte uma solicitação mecânica em informação sobre a ocorrência de um evento.

O alerta não demonstra presença ou ausência de trinca, não torna a cerâmica visível e não substitui a contenção. O limiar no ponto de montagem também não é automaticamente o limite de aceleração na peça. A especificação do monitoramento abrange tecnologia, posição, calibração, reposição ou rearme, vida e custo. O registro de impacto é a opção OP-01 da equipe, associada à M3 nesta comparação e fora das 21 elementares da M1 selecionada.

\begin{landscape}
\begin{figure}[H]
\centering\begin{tikzpicture}[x=1mm,y=1mm]
\tikzset{
 fn/.style={draw=cinzamedio,fill=white,align=center,font=\footnotesize,inner xsep=1mm,inner ysep=1mm,text width=22mm,minimum height=13mm,execute at begin node=\setstretch{1}},
 fp/.style={fn,fill=cinzaclaro,minimum height=14mm},
 material/.style={-{Stealth[length=1.7mm]},draw=cinzaescuro,line width=.8pt},
 energia/.style={-{Stealth[length=1.5mm]},draw=cinzamedio,dashed,line width=.6pt},
 sinal/.style={-{Stealth[length=1.5mm]},draw=cinzamedio,dash pattern=on 3pt off 2pt on .6pt off 2pt,line width=.6pt},
 rel/.style={-{Stealth[length=1mm]},draw=cinzamedio,line width=.45pt},
 porta/.style={align=left,font=\footnotesize,inner sep=0,execute at begin node=\setstretch{1}},
 legenda/.style={font=\footnotesize,anchor=west,execute at begin node=\setstretch{1}}
}

\fill[cinzaclaro,opacity=.45] (26,8) rectangle (52,-75);
\node[fp] (p0) at (39,0) {{\scriptsize P4}\\\textbf{Acomodar louça}};
\node[fn] (e00) at (39,-19) {{\scriptsize E10}\\\textbf{Receber louça}};
\node[fn] (e01) at (39,-34) {{\scriptsize E11}\\\textbf{Ajustar apoios}};
\node[fn] (e02) at (39,-49) {{\scriptsize E12}\\\textbf{Posicionar louça}};
\draw[energia] (39,-82)--(39,-75);
\node[porta,anchor=north,align=center,text width=25mm,font=\scriptsize] at (39,-84) {Trabalho do operador};
\fill[cinzaclaro,opacity=.45] (53,8) rectangle (79,-75);
\node[fp] (p1) at (66,0) {{\scriptsize P1}\\\textbf{Imobilizar louça}};
\node[fn] (e10) at (66,-19) {{\scriptsize E01}\\\textbf{Travar apoios}};
\node[fn] (e11) at (66,-34) {{\scriptsize E02}\\\textbf{Distribuir pressão}};
\node[fn] (e12) at (66,-49) {{\scriptsize E03}\\\textbf{Restringir movimento}};
\draw[energia] (66,-82)--(66,-75);
\node[porta,anchor=north,align=center,text width=25mm,font=\scriptsize] at (66,-84) {Trabalho do operador};
\fill[cinzaclaro,opacity=.45] (80,8) rectangle (106,-75);
\node[fp] (p2) at (93,0) {{\scriptsize P2}\\\textbf{Atenuar esforços}};
\node[fn] (e20) at (93,-19) {{\scriptsize E04}\\\textbf{Absorver impacto}};
\node[fn] (e21) at (93,-34) {{\scriptsize E05}\\\textbf{Isolar vibração}};
\node[fn] (e22) at (93,-49) {{\scriptsize E06}\\\textbf{Desviar compressão}};
\draw[energia] (93,-82)--(93,-75);
\node[porta,anchor=north,align=center,text width=25mm,font=\scriptsize] at (93,-84) {Choque, vibração e compressão};
\fill[cinzaclaro,opacity=.45] (107,8) rectangle (133,-75);
\node[fp] (p3) at (120,0) {{\scriptsize P3}\\\textbf{Conter fragmentos}};
\node[fn] (e30) at (120,-19) {{\scriptsize E07}\\\textbf{Envolver louça}};
\node[fn] (e31) at (120,-34) {{\scriptsize E08}\\\textbf{Reter fragmentos}};
\node[fn] (e32) at (120,-49) {{\scriptsize E09}\\\textbf{Vedar juntas}};
\draw[energia] (120,-82)--(120,-75);
\node[porta,anchor=north,align=center,text width=25mm,font=\scriptsize] at (120,-84) {Energia da quebra};
\fill[cinzaclaro,opacity=.45] (134,8) rectangle (160,-75);
\node[fp] (p4) at (147,0) {{\scriptsize P5}\\\textbf{Movimentar conjunto}};
\node[fn] (e40) at (147,-19) {{\scriptsize E13}\\\textbf{Receber força}};
\node[fn] (e41) at (147,-34) {{\scriptsize E14}\\\textbf{Transmitir força}};
\node[fn] (e42) at (147,-49) {{\scriptsize E15}\\\textbf{Estabilizar conjunto}};
\draw[energia] (147,-82)--(147,-75);
\node[porta,anchor=north,align=center,text width=25mm,font=\scriptsize] at (147,-84) {Trabalho do operador};
\fill[cinzaclaro,opacity=.45] (161,8) rectangle (187,-75);
\node[fp] (p5) at (174,0) {{\scriptsize P6}\\\textbf{Informar condição}};
\node[fn] (e50) at (174,-19) {{\scriptsize E16}\\\textbf{Expor superfície}};
\node[fn] (e51) at (174,-34) {{\scriptsize E17}\\\textbf{Indicar fechamento}};
\node[fn] (e52) at (174,-49) {{\scriptsize E18}\\\textbf{Exibir identidade}};
\fill[cinzaclaro,opacity=.45] (188,8) rectangle (214,-75);
\node[fp] (p6) at (201,0) {{\scriptsize P7}\\\textbf{Liberar louça}};
\node[fn] (e60) at (201,-19) {{\scriptsize E19}\\\textbf{Destravar apoios}};
\node[fn] (e61) at (201,-34) {{\scriptsize E20}\\\textbf{Abrir barreira}};
\node[fn] (e62) at (201,-49) {{\scriptsize E21}\\\textbf{Remover louça}};
\draw[energia] (201,-82)--(201,-75);
\node[porta,anchor=north,align=center,text width=25mm,font=\scriptsize] at (201,-84) {Trabalho do operador};
\draw[material] (p0.east) -- (p1.west);
\draw[material] (p1.east) -- (p2.west);
\draw[material] (p2.east) -- (p3.west);
\draw[material] (p3.east) -- (p4.west);
\draw[material] (p4.east) -- (p5.west);
\draw[material] (p5.east) -- (p6.west);
\draw[material] (0,0)--(p0.west);
\draw[material] (p6.east)--(244,0);
\draw[material] (214,-27)--(244,-27);
\node[porta,anchor=south west,text width=24mm] at (0,2) {Louça\\desembalada};
\node[porta,anchor=south west,text width=29mm] at (216,2) {Louça íntegra\\entregue};
\node[porta,anchor=south west,text width=29mm] at (216,-25) {Fragmentos retidos, se houver quebra};
\draw[energia] (93,8)--(93,17);
\node[porta,anchor=south,align=center,text width=48mm] at (93,19) {Energia dissipada};
\draw[sinal] (174,8)--(174,17);
\node[porta,anchor=south,align=center,text width=64mm] at (174,19) {Condição, fechamento e identidade};
\draw[material] (18,-104)--(32,-104);\node[legenda] at (35,-104) {Material};
\draw[energia] (88,-104)--(102,-104);\node[legenda] at (105,-104) {Energia};
\draw[sinal] (158,-104)--(172,-104);\node[legenda] at (175,-104) {Sinal};
\end{tikzpicture}

\caption{Modelagem M1: proteção direta da louça descoberta}
\label{fig:modelagem-m1}
\fonte{Elaboração da equipe, conforme o catálogo e a estrutura funcional vigentes.}
\notaquadro{As parciais organizam os fluxos; os blocos inferiores mostram suas elementares. As sete complementares são comuns às três modelagens.}
\end{figure}
\end{landscape}

\begin{landscape}
\begin{figure}[H]
\centering\begin{tikzpicture}[x=1mm,y=1mm]
\tikzset{
 fn/.style={draw=cinzamedio,fill=white,align=center,font=\footnotesize,inner xsep=1mm,inner ysep=1mm,text width=22mm,minimum height=13mm,execute at begin node=\setstretch{1}},
 fp/.style={fn,fill=cinzaclaro,minimum height=14mm},
 material/.style={-{Stealth[length=1.7mm]},draw=cinzaescuro,line width=.8pt},
 energia/.style={-{Stealth[length=1.5mm]},draw=cinzamedio,dashed,line width=.6pt},
 sinal/.style={-{Stealth[length=1.5mm]},draw=cinzamedio,dash pattern=on 3pt off 2pt on .6pt off 2pt,line width=.6pt},
 rel/.style={-{Stealth[length=1mm]},draw=cinzamedio,line width=.45pt},
 porta/.style={align=left,font=\footnotesize,inner sep=0,execute at begin node=\setstretch{1}},
 legenda/.style={font=\footnotesize,anchor=west,execute at begin node=\setstretch{1}}
}

\fill[cinzaclaro,opacity=.45] (26,8) rectangle (52,-75);
\node[fp] (p0) at (39,0) {{\scriptsize P4}\\\textbf{Acomodar volume}};
\node[fn] (e00) at (39,-19) {{\scriptsize E10}\\\textbf{Receber volume}};
\node[fn] (e01) at (39,-34) {{\scriptsize E11}\\\textbf{Ajustar apoios}};
\node[fn] (e02) at (39,-49) {{\scriptsize E12}\\\textbf{Posicionar volume}};
\draw[energia] (39,-82)--(39,-75);
\node[porta,anchor=north,align=center,text width=25mm,font=\scriptsize] at (39,-84) {Trabalho do operador};
\fill[cinzaclaro,opacity=.45] (53,8) rectangle (79,-75);
\node[fp] (p1) at (66,0) {{\scriptsize P1}\\\textbf{Imobilizar volume}};
\node[fn] (e10) at (66,-19) {{\scriptsize E01}\\\textbf{Travar apoios}};
\node[fn] (e11) at (66,-34) {{\scriptsize E02}\\\textbf{Distribuir pressão}};
\node[fn] (e12) at (66,-49) {{\scriptsize E03}\\\textbf{Restringir movimento}};
\draw[energia] (66,-82)--(66,-75);
\node[porta,anchor=north,align=center,text width=25mm,font=\scriptsize] at (66,-84) {Trabalho do operador};
\fill[cinzaclaro,opacity=.45] (80,8) rectangle (106,-75);
\node[fp] (p2) at (93,0) {{\scriptsize P2}\\\textbf{Atenuar esforços}};
\node[fn] (e20) at (93,-19) {{\scriptsize E04}\\\textbf{Absorver impacto}};
\node[fn] (e21) at (93,-34) {{\scriptsize E05}\\\textbf{Isolar vibração}};
\node[fn] (e22) at (93,-49) {{\scriptsize E06}\\\textbf{Desviar compressão}};
\draw[energia] (93,-82)--(93,-75);
\node[porta,anchor=north,align=center,text width=25mm,font=\scriptsize] at (93,-84) {Choque, vibração e compressão};
\fill[cinzaclaro,opacity=.45] (107,8) rectangle (133,-75);
\node[fp] (p3) at (120,0) {{\scriptsize P3}\\\textbf{Conter fragmentos}};
\node[fn] (e30) at (120,-19) {{\scriptsize E07}\\\textbf{Envolver volume}};
\node[fn] (e31) at (120,-34) {{\scriptsize E08}\\\textbf{Reter fragmentos}};
\node[fn] (e32) at (120,-49) {{\scriptsize E09}\\\textbf{Vedar juntas}};
\draw[energia] (120,-82)--(120,-75);
\node[porta,anchor=north,align=center,text width=25mm,font=\scriptsize] at (120,-84) {Energia da quebra};
\fill[cinzaclaro,opacity=.45] (134,8) rectangle (160,-75);
\node[fp] (p4) at (147,0) {{\scriptsize P5}\\\textbf{Movimentar conjunto}};
\node[fn] (e40) at (147,-19) {{\scriptsize E13}\\\textbf{Receber força}};
\node[fn] (e41) at (147,-34) {{\scriptsize E14}\\\textbf{Transmitir força}};
\node[fn] (e42) at (147,-49) {{\scriptsize E15}\\\textbf{Estabilizar conjunto}};
\draw[energia] (147,-82)--(147,-75);
\node[porta,anchor=north,align=center,text width=25mm,font=\scriptsize] at (147,-84) {Trabalho do operador};
\fill[cinzaclaro,opacity=.45] (161,8) rectangle (187,-75);
\node[fp] (p5) at (174,0) {{\scriptsize P6}\\\textbf{Informar condição}};
\node[fn] (e50) at (174,-19) {{\scriptsize A01}\\\textbf{Abrir acesso}};
\node[fn] (e51) at (174,-34) {{\scriptsize E16}\\\textbf{Expor superfície}};
\node[fn] (e52) at (174,-49) {{\scriptsize E17}\\\textbf{Indicar fechamento}};
\node[fn] (e53) at (174,-64) {{\scriptsize E18}\\\textbf{Exibir identidade}};
\draw[energia] (174,-82)--(174,-75);
\node[porta,anchor=north,align=center,text width=25mm,font=\scriptsize] at (174,-84) {Trabalho do operador};
\fill[cinzaclaro,opacity=.45] (188,8) rectangle (214,-75);
\node[fp] (p6) at (201,0) {{\scriptsize P7}\\\textbf{Liberar volume}};
\node[fn] (e60) at (201,-19) {{\scriptsize E19}\\\textbf{Destravar apoios}};
\node[fn] (e61) at (201,-34) {{\scriptsize E20}\\\textbf{Abrir barreira}};
\node[fn] (e62) at (201,-49) {{\scriptsize E21}\\\textbf{Remover volume}};
\draw[energia] (201,-82)--(201,-75);
\node[porta,anchor=north,align=center,text width=25mm,font=\scriptsize] at (201,-84) {Trabalho do operador};
\draw[material] (p0.east) -- (p1.west);
\draw[material] (p1.east) -- (p2.west);
\draw[material] (p2.east) -- (p3.west);
\draw[material] (p3.east) -- (p4.west);
\draw[material] (p4.east) -- (p5.west);
\draw[material] (p5.east) -- (p6.west);
\draw[material] (0,0)--(p0.west);
\draw[material] (p6.east)--(244,0);
\draw[material] (214,-27)--(244,-27);
\node[porta,anchor=south west,text width=24mm] at (0,2) {Louça na\\embalagem\\original};
\node[porta,anchor=south west,text width=29mm] at (216,2) {Volume\\embalado\\entregue};
\node[porta,anchor=south west,text width=29mm] at (216,-25) {Fragmentos retidos, se houver quebra};
\draw[energia] (93,8)--(93,17);
\node[porta,anchor=south,align=center,text width=48mm] at (93,19) {Energia dissipada};
\draw[sinal] (174,8)--(174,17);
\node[porta,anchor=south,align=center,text width=64mm] at (174,19) {Condição, fechamento e identidade};
\draw[material] (18,-104)--(32,-104);\node[legenda] at (35,-104) {Material};
\draw[energia] (88,-104)--(102,-104);\node[legenda] at (105,-104) {Energia};
\draw[sinal] (158,-104)--(172,-104);\node[legenda] at (175,-104) {Sinal};
\end{tikzpicture}

\caption{Modelagem M2: proteção do volume na embalagem original}
\label{fig:modelagem-m2}
\fonte{Elaboração da equipe, conforme o catálogo e a estrutura funcional vigentes.}
\notaquadro{A01, Abrir acesso, é adicional de M2/M3. O acesso pela embalagem interna deve preservar a condição de inspeção com o conjunto fechado.}
\end{figure}
\end{landscape}

\begin{landscape}
\begin{figure}[H]
\centering\begin{tikzpicture}[x=1mm,y=.88mm]
\tikzset{
 fn/.style={draw=cinzamedio,fill=white,align=center,font=\footnotesize,inner xsep=1mm,inner ysep=1mm,text width=22mm,minimum height=13mm,execute at begin node=\setstretch{1}},
 fp/.style={fn,fill=cinzaclaro,minimum height=14mm},
 material/.style={-{Stealth[length=1.7mm]},draw=cinzaescuro,line width=.8pt},
 energia/.style={-{Stealth[length=1.5mm]},draw=cinzamedio,dashed,line width=.6pt},
 sinal/.style={-{Stealth[length=1.5mm]},draw=cinzamedio,dash pattern=on 3pt off 2pt on .6pt off 2pt,line width=.6pt},
 rel/.style={-{Stealth[length=1mm]},draw=cinzamedio,line width=.45pt},
 porta/.style={align=left,font=\footnotesize,inner sep=0,execute at begin node=\setstretch{1}},
 legenda/.style={font=\footnotesize,anchor=west,execute at begin node=\setstretch{1}}
}

\tikzset{fn/.append style={minimum height=11.5mm},fp/.append style={minimum height=12.5mm}}
\fill[cinzaclaro,opacity=.45] (26,8) rectangle (52,-75);
\node[fp] (p0) at (39,0) {{\scriptsize P4}\\\textbf{Acomodar volume}};
\node[fn] (e00) at (39,-19) {{\scriptsize E10}\\\textbf{Receber volume}};
\node[fn] (e01) at (39,-34) {{\scriptsize E11}\\\textbf{Ajustar apoios}};
\node[fn] (e02) at (39,-49) {{\scriptsize E12}\\\textbf{Posicionar volume}};
\draw[energia] (39,-82)--(39,-75);
\node[porta,anchor=north,align=center,text width=25mm,font=\scriptsize] at (39,-84) {Trabalho do operador};
\fill[cinzaclaro,opacity=.45] (53,8) rectangle (79,-75);
\node[fp] (p1) at (66,0) {{\scriptsize P1}\\\textbf{Imobilizar volume}};
\node[fn] (e10) at (66,-19) {{\scriptsize E01}\\\textbf{Travar apoios}};
\node[fn] (e11) at (66,-34) {{\scriptsize E02}\\\textbf{Distribuir pressão}};
\node[fn] (e12) at (66,-49) {{\scriptsize E03}\\\textbf{Restringir movimento}};
\draw[energia] (66,-82)--(66,-75);
\node[porta,anchor=north,align=center,text width=25mm,font=\scriptsize] at (66,-84) {Trabalho do operador};
\fill[cinzaclaro,opacity=.45] (80,8) rectangle (106,-75);
\node[fp] (p2) at (93,0) {{\scriptsize P2}\\\textbf{Atenuar esforços}};
\node[fn] (e20) at (93,-19) {{\scriptsize E04}\\\textbf{Absorver impacto}};
\node[fn] (e21) at (93,-34) {{\scriptsize E05}\\\textbf{Isolar vibração}};
\node[fn] (e22) at (93,-49) {{\scriptsize E06}\\\textbf{Desviar compressão}};
\draw[energia] (93,-82)--(93,-75);
\node[porta,anchor=north,align=center,text width=25mm,font=\scriptsize] at (93,-84) {Choque, vibração e compressão};
\fill[cinzaclaro,opacity=.45] (107,8) rectangle (133,-75);
\node[fp] (p3) at (120,0) {{\scriptsize P3}\\\textbf{Conter fragmentos}};
\node[fn] (e30) at (120,-19) {{\scriptsize E07}\\\textbf{Envolver volume}};
\node[fn] (e31) at (120,-34) {{\scriptsize E08}\\\textbf{Reter fragmentos}};
\node[fn] (e32) at (120,-49) {{\scriptsize E09}\\\textbf{Vedar juntas}};
\draw[energia] (120,-82)--(120,-75);
\node[porta,anchor=north,align=center,text width=25mm,font=\scriptsize] at (120,-84) {Energia da quebra};
\fill[cinzaclaro,opacity=.45] (134,8) rectangle (160,-75);
\node[fp] (p4) at (147,0) {{\scriptsize P5}\\\textbf{Movimentar conjunto}};
\node[fn] (e40) at (147,-19) {{\scriptsize E13}\\\textbf{Receber força}};
\node[fn] (e41) at (147,-34) {{\scriptsize E14}\\\textbf{Transmitir força}};
\node[fn] (e42) at (147,-49) {{\scriptsize E15}\\\textbf{Estabilizar conjunto}};
\draw[energia] (147,-82)--(147,-75);
\node[porta,anchor=north,align=center,text width=25mm,font=\scriptsize] at (147,-84) {Trabalho do operador};
\fill[cinzaclaro,opacity=.45] (161,8) rectangle (187,-75);
\node[fp] (p5) at (174,0) {{\scriptsize P6}\\\textbf{Informar condição}};
\node[fn] (e50) at (174,-19) {{\scriptsize A01}\\\textbf{Abrir acesso}};
\node[fn] (e51) at (174,-34) {{\scriptsize E16}\\\textbf{Expor superfície}};
\node[fn] (e52) at (174,-49) {{\scriptsize E17}\\\textbf{Indicar fechamento}};
\node[fn] (e53) at (174,-64) {{\scriptsize E18}\\\textbf{Exibir identidade}};
\draw[energia] (174,-82)--(174,-75);
\node[porta,anchor=north,align=center,text width=25mm,font=\scriptsize] at (174,-84) {Trabalho do operador};
\fill[cinzaclaro,opacity=.45] (188,8) rectangle (214,-75);
\node[fp] (p6) at (201,0) {{\scriptsize P7}\\\textbf{Liberar volume}};
\node[fn] (e60) at (201,-19) {{\scriptsize E19}\\\textbf{Destravar apoios}};
\node[fn] (e61) at (201,-34) {{\scriptsize E20}\\\textbf{Abrir barreira}};
\node[fn] (e62) at (201,-49) {{\scriptsize E21}\\\textbf{Remover volume}};
\draw[energia] (201,-82)--(201,-75);
\node[porta,anchor=north,align=center,text width=25mm,font=\scriptsize] at (201,-84) {Trabalho do operador};
\draw[material] (p0.east) -- (p1.west);
\draw[material] (p1.east) -- (p2.west);
\draw[material] (p2.east) -- (p3.west);
\draw[material] (p3.east) -- (p4.west);
\draw[material] (p4.east) -- (p5.west);
\draw[material] (p5.east) -- (p6.west);
\draw[material] (0,0)--(p0.west);
\draw[material] (p6.east)--(244,0);
\draw[material] (214,-27)--(244,-27);
\node[porta,anchor=south west,text width=24mm] at (0,2) {Louça na\\embalagem\\original};
\node[porta,anchor=south west,text width=29mm] at (216,2) {Volume\\embalado\\entregue};
\node[porta,anchor=south west,text width=29mm] at (216,-25) {Fragmentos retidos, se houver quebra};
\draw[energia] (93,8)--(93,17);
\node[porta,anchor=south,align=center,text width=48mm] at (93,19) {Energia dissipada};
\draw[sinal] (174,8)--(174,17);
\node[porta,anchor=south,align=center,text width=64mm] at (174,19) {Condição, fechamento e identidade};
\node[fp,text width=38mm] (monitor) at (27,-107) {\textbf{Monitorar impacto}};
\node[fn,text width=30mm] (o0) at (76,-107) {{\scriptsize O01}\\\textbf{Detectar aceleração}};
\node[fn,text width=30mm] (o1) at (121,-107) {{\scriptsize O02}\\\textbf{Comparar limiar}};
\node[fn,text width=30mm] (o2) at (166,-107) {{\scriptsize O03}\\\textbf{Registrar evento}};
\node[fn,text width=30mm] (o3) at (211,-107) {{\scriptsize O04}\\\textbf{Exibir alerta}};
\node[porta,anchor=south west,text width=40mm] at (0,-97) {Choque no volume};
\draw[energia] (0,-107)--(monitor.west);
\draw[energia] (monitor.east)--(o0.west);
\draw[sinal] (o0.east)--(o1.west);
\draw[sinal] (o1.east)--(o2.west);
\draw[sinal] (o2.east)--(o3.west);
\draw[sinal] (o3.east)--(244,-107);
\node[porta,anchor=north east,text width=40mm,align=right] at (244,-117) {Alerta de evento};
\draw[material] (18,-131)--(32,-131);\node[legenda] at (35,-131) {Material};
\draw[energia] (88,-131)--(102,-131);\node[legenda] at (105,-131) {Energia};
\draw[sinal] (158,-131)--(172,-131);\node[legenda] at (175,-131) {Sinal};
\end{tikzpicture}

\caption{Modelagem M3: cadeia principal da M2 com ramo de monitoramento}
\label{fig:modelagem-m3}
\fonte{Elaboração da equipe, conforme o catálogo e a estrutura funcional vigentes.}
\notaquadro{Monitorar impacto opera em paralelo à cadeia material da M2. O alerta de evento não substitui Expor superfície nem comprova integridade.}
\end{figure}
\end{landscape}



\subsubsection{Seleção da Modelagem}

A seleção usa a matriz de Pugh para comparar estruturas no mesmo nível de abstração. As notas expressam julgamentos relativos à referência, apoiados nas diferenças de fronteira e nas hipóteses de solução disponíveis. A pontuação orienta o desenvolvimento; a comprovação de desempenho depende das soluções físicas e de suas verificações \citep{back2008,mitpugh2015}.

M2 é a referência, ou \textit{datum}, e a escala de comparação é $-1$, 0 e $+1$. A nota positiva indica vantagem prevista frente à referência; a negativa, desvantagem; zero registra equivalência adotada ou ausência de vantagem comparativa demonstrada. Os pesos são os graus de importância do cliente, de 0 a 5, derivados do Mudge/Pareto. Eles não são os pontos de importância técnica dos RP no QFD. A comparação contempla os 23 requisitos, inclusive os que restringem funções e soluções, conforme a Tabela~\ref{tab:pugh-funcional}.

\begin{table}[H]
\caption{Primeira Pugh: comparação ponderada das modelagens funcionais}
\label{tab:pugh-funcional}
\centering\footnotesize
\begin{tabularx}{\textwidth}{@{}L{0.55cm}L{0.85cm}YL{0.75cm}*{3}{>{\centering\arraybackslash}p{0.65cm}}@{}}
\toprule\textbf{RC}&\textbf{RP}&\textbf{Requisito do cliente}&\textbf{Peso}&\textbf{M1}&\textbf{M2}&\textbf{M3}\\\midrule
\linhasquadro{quadros/pugh-modelagens-atuais}
\midrule
\multicolumn{4}{l}{\textbf{Total ponderado}}&\textbf{+7}&\textbf{0}&\textbf{$-4$}\\
\bottomrule\end{tabularx}
\fonte{Avaliação funcional dos 23 requisitos; pesos derivados do Mudge/Pareto.}
\notaquadro{M2 é a referência. Zero não comprova cumprimento das metas. RP19 continua obrigatório mesmo com peso zero no Pareto.}
\end{table}

O total ponderado de cada modelagem é dado por $S_j=\sum_{i=1}^{23}w_i n_{ij}$, em que $w_i$ é o grau de importância do requisito e $n_{ij}$ é a nota relativa à M2. M1 apresenta cinco vantagens, quatro desvantagens e quatorze zeros; M3 apresenta uma vantagem, cinco desvantagens e dezessete zeros. Os 23 zeros da M2 decorrem de sua condição de referência e não de desempenho aprovado.

\noindent\textbf{Fundamentação das diferenças.} As notas da Tabela~\ref{tab:fundamentacao-pugh} relacionam cada diferença à fronteira material e às capacidades funcionais. Elas descrevem potenciais e contrapartidas a verificar no mesmo mix, percurso e condições. A retirada da caixa muda a interface e o preparo; a inclusão de monitoramento muda a informação e as tarefas de verificação. Nenhuma dessas mudanças demonstra sozinha melhor proteção física.

\begin{table}[H]
\caption{Fundamentação das notas diferentes de zero}
\label{tab:fundamentacao-pugh}
\centering\footnotesize
\begin{tabularx}{\textwidth}{L{1.5cm}*{2}{>{\centering\arraybackslash}p{0.7cm}}Y}
\toprule\textbf{RC / RP}&\textbf{M1}&\textbf{M3}&\textbf{Diferença e condição da avaliação}\\\midrule
B / RP2&+1&0&Na M1 a restrição age diretamente na peça; M2/M3 mantêm a folga interna da caixa. O julgamento depende de contatos que limitem pressão e deslocamento sem danificar a peça.\\
C / RP4&+1&0&M1 retira o papelão opaco do caminho visual. A barreira do dispositivo permanece e deve permitir conferir a área crítica com o conjunto fechado.\\
H / RP6&+1&+1&M1 reduz a necessidade de abrir a caixa com estilete na obra; M3 acrescenta aviso prévio de evento. São hipóteses de retirada; o alerta não comprova retenção nem abertura segura.\\
D / RP7&$-1$&0&M1 exige regular apoios para geometrias irregulares de cubas e bacias. M2/M3 atuam sobre o volume externo da caixa. Comparar a mesma faixa de peças.\\
V / RP23&0&$-1$&M3 acrescenta aquisição e reposição/rearme do monitoramento. M1 recebe equivalência comparativa de custo; o critério econômico inclui preparo, retorno e vida de serviço.\\
I / RP10&$-1$&$-1$&M1 acrescenta desembalagem e transferência no preparo; M3 exige interpretar o alerta. Confirmar ações sob uma fronteira operacional comum.\\
J / RP14&+1&0&Eliminar a caixa interna pode reduzir o envelope. Medir o conjunto carregado por família; não presumir economia de frete.\\
P / RP13&+1&0&M1 reduz a dependência do papelão interno. A vantagem decorre dessa fronteira; resistência úmida é critério do conjunto.\\
O / RP18&0&$-1$&A nota da M3 considera reposição de indicador. Uma tecnologia reutilizável pode alterar o julgamento; a vida do conjunto precisa de ensaio.\\
U / RP22&$-1$&$-1$&M1 exige regulagem dos apoios à peça; M3 acrescenta instalação/rearme do monitoramento. Confirmar a contagem de ações.\\
Q / RP15&$-1$&$-1$&M1 exige cuidado com limpeza e desgaste dos contatos; M3 acrescenta verificação e manutenção do monitoramento conforme a tecnologia.\\
\bottomrule\end{tabularx}
\fonte{Análise comparativa das interfaces e dos efeitos das três modelagens.}
\end{table}

A comparação adota folga interna de aproximadamente 20\,mm por face como premissa geométrica da embalagem. Essa grandeza é distinta do deslocamento residual do RP2. O efeito da umidade sobre o papelão tem suporte na pesquisa de \citet{brown2019}, sem que um fator único de perda seja transferido a todas as caixas. As metas da Tabela~\ref{tab:especificacoes-meta} permanecem os critérios de aceitação do produto.

Os zeros em proteção, movimentação, pega, compatibilidade com o veículo, estabilidade, empilhamento e uniformidade registram que esta comparação não estabeleceu uma diferença que justificasse outra nota. Orientação visual, desmontagem, separação, modularidade e identidade também permanecem capacidades exigidas das três estruturas. O custo da M1 recebeu zero sem cotação comparável, não porque se tenha demonstrado custo igual. A equivalência é de comparação funcional, sem substituir a comprovação das respectivas metas.

\noindent\textbf{Sensibilidade da comparação.} Retirar a vantagem de M1 em imobilização reduz seu total para +2; retirar também a vantagem em inspeção reduz para $-1$. A preferência depende, portanto, da preservação desses benefícios no mesmo conjunto. A análise de sensibilidade não altera as notas adotadas na matriz.

\noindent\textbf{Modelagem selecionada: M1, proteção direta da peça descoberta.} A M1 obtém +7, frente a 0 da M2 e $-4$ da M3. A restrição atua diretamente sobre a louça e a retirada da caixa original elimina o papelão opaco do caminho de observação. Essas diferenças fundamentam as vantagens previstas em imobilização e inspeção, os critérios de maior influência na seleção. A redução potencial do envelope e da dependência do papelão úmido também contribui para o resultado.

A escolha considera as contrapartidas de desembalagem e transferência no distribuidor, ajuste para geometrias irregulares e limpeza dos contatos. O atendimento exige pressão de apoio sem dano, conferência da área crítica com a barreira fechada e contenção durante a retirada. M2 permanece a referência de contingência; o alerta adicional de M3 não substitui inspeção nem compensa, nesta comparação, suas desvantagens de custo, uso e manutenção. Com esses critérios e premissas, a estrutura funcional selecionada é a M1.

